\scanpage{77}
- \mspage{209} -
\margin{n$^\circ$ 391}
(VII) On a déterminé ci-dessus les racines positives ; leur somme est donc :
\[
2\rho=5\varepsilon_1+3\varepsilon_2+\varepsilon_3+9\varepsilon_4
=16\alpha_1+30\alpha_2+42\alpha_3+22\alpha_4.
\]

(VIII) On vient de voir que l'on a $\Gamma(R)=\Gamma_o(R)$, l'indice de connexion est donc $1$.

(IX) La famille des exposants a $4$ termes, et puisque $h=12$, les entiers $1$, $5$, $7$ et $11$ étrangers à $12$, doivent figurer dans cette famille, (\S\ 8, n$^\circ$ 5, prop. 5); ce sont par suite tous les exposants de $W(R)$.

(X) On déduit de (IX) et du \S\ 8, n$^\circ$ 3, cor. 1 de la prop. 3 que l'ordre du groupe de Weyl $W(R)$ est :
\[
2.6.8.12=2^7.3^2=1152.
\]

(XI) Le même raisonnement qu'au n$^\circ$ 6 montre que le seul automorphisme du schéma de Dynkin $F_4$ est l'identité.

10. \uline{Système de type} $E_8$.

(I) Considérons dans $V=R^8$ le réseau $L_3$, (n$^\circ$ 4) ; l'ensemble $R$ des vecteurs $u$ de $L_3$ tels que $(u|u)=2$ est formé des vecteurs :
\[
\pm\varepsilon_i\pm\varepsilon_j,\ (i\ne j),\quad
\tfrac12\sum_{i=1}^{8}(-1)^{\unsure{\vartheta|\delta}(i)}\varepsilon_i
\]
où, (en vertu de la définition de $L_3$), la somme des $\unsure{\vartheta|\delta}(i)$ doit être \uline{paire} ; cela résulte de nouveau du fait qu'un vecteur de $R$ ne peut avoir de composante de valeur absolue $3/2$. On a vu, (n$^\circ$ 4), que $(u|v)\in Z$ quels que soient $u$, $v$ dans $R$, et le même raisonnement qu'au n$^\circ$ 5 montre que $R$ est un \uline{système de racines réduit}, ayant $n=\binom82.4+2^7=240$ éléments.

(II) Munissons $R^8$ de l'ordre \uline{opposé} à l'\uline{ordre lexicographique} défini par les $\varepsilon_i$, de sorte que pour cet ordre on a $\varepsilon_i<\varepsilon_{i+1}$ pour $1\leq i\leq7$. Considérons les $8$ vecteurs :
\mspage{76}
